\chapter{Carga y visualización de datos}
\label{ch:carga}

Con la calibración fijada (capítulo~\ref{ch:calibracion}), cargar un espectro es
inmediato: Fitbauer lo dobla, construye el eje de velocidad y lo muestra
normalizado. Este capítulo detalla los formatos, el \emph{folding}, la
normalización y las opciones de visualización.

\section{Formatos admitidos}

\begin{center}
\begin{tabular}{@{}lp{0.62\textwidth}@{}}
\toprule
\textbf{Extensión} & \textbf{Contenido} \\
\midrule
\ruta{.ws5}, \ruta{.adt} & \textbf{Cuentas en bruto} por canal (típ.\ 512
canales). No traen eje de velocidad: se doblan y se calibran con la \param{Vmax}
fijada. \\[4pt]
\ruta{.csv}, \ruta{.txt}, \ruta{.dat}, \ruta{.exp} & Espectro \textbf{ya doblado}
en espacio de velocidad (columnas velocidad/cuentas). Traen su propio eje: no se
doblan y \emph{sustituyen} la calibración (§\ref{sec:calib-fija}). \\
\bottomrule
\end{tabular}
\end{center}

\begin{nota}
Los ficheros \ruta{.ws5}/\ruta{.adt} son listas de enteros (una cuenta por
canal). Los de velocidad (\ruta{.csv}/\ruta{.txt}/…) son dos columnas
—velocidad y cuentas/transmisión— separadas por espacios, tabuladores o comas;
Fitbauer tolera cabeceras y líneas de comentario. Si un fichero de velocidad no
carga, revise que tenga exactamente dos columnas numéricas por línea.
\end{nota}

\section{Cargar un espectro}

\menu{Archivo > Cargar…} abre el selector de ficheros. Al elegir un espectro de
cuentas en bruto, Fitbauer:

\begin{enumerate}[nosep]
  \item Lee las cuentas por canal.
  \item Detecta el \textbf{centro de folding} (o lo lee de un fichero
        \ruta{.RES} de NORMOS si existe junto al espectro).
  \item \textbf{Dobla} las dos ramas del espectro triangular y \textbf{recorta}
        unos pocos canales de cada extremo (§\ref{sec:folding} explica por qué).
  \item \textbf{Normaliza} a línea base $\approx 1$ y estima el ruido de Poisson.
  \item Construye el eje de velocidad con la \param{Vmax} de la calibración
        fijada.
\end{enumerate}

\autofig[\textwidth]{fig_load_alphafe}{Un mismo $\alpha$-Fe \textbf{sin doblar}
(izq., 512 canales) y \textbf{doblado} (dcha., en velocidad). Sin doblar, el
transductor mide el sexteto \emph{dos veces} —una en cada rama del barrido—,
espejadas respecto al \textbf{centro de folding} (línea roja). El folding suma
esas dos mitades: reduce el ruido y produce el espectro simétrico de la derecha.}

\begin{nota}
La barra de estado y la etiqueta del fichero informan del número de canales, el
centro de folding detectado y el factor de normalización. \menu{Archivo > Abrir
recientes} reabre los últimos espectros usados.
\end{nota}

\section{El folding y el centro}
\label{sec:folding}

El transductor Mössbauer recorre un ciclo de velocidad \textbf{triangular}: cada
velocidad se mide dos veces, una en la rama ascendente y otra en la descendente.
El \emph{folding} \textbf{suma esas dos mitades} para promediarlas: mejora la
relación señal/ruido y produce un único espectro simétrico. El punto clave es el
\textbf{centro}: el canal (posiblemente fraccionario) donde se pliega.

\begin{itemize}[nosep]
  \item Fitbauer detecta el centro automáticamente (o lo lee de un \ruta{.RES}).
  \item El control \param{Centro} del panel de calibración lo ajusta a mano: el
        espectro se \textbf{re-dobla en vivo} al cambiarlo.
  \item \menu{Ajuste > Buscar centro} lo refina por minimización.
  \item La casilla \param{Ajustar folding point dentro del ajuste} lo incluye
        como \textbf{parámetro libre del ajuste discreto}: en cada iteración el
        espectro se re-dobla con el centro propuesto, se renormaliza y se
        recalculan los pesos estadísticos; al terminar, el valor ajustado se
        vuelca al control \param{Centro} y los datos quedan doblados con él.
        Útil cuando el residuo muestra pares antisimétricos alrededor de las
        líneas; conviene partir de un valor razonable y con pocos parámetros
        libres.
\end{itemize}

\paragraph{Cómo dobla Fitbauer.} Tres detalles separan un doblado cuidadoso de
uno descuidado, y los tres se notan en la anchura de línea recuperada.

\begin{itemize}[nosep]
  \item \textbf{Interpolación cúbica.} El punto de simetría casi nunca cae en un
        canal entero, así que hay que interpolar el canal emparejado. Fitbauer
        usa un B-spline cúbico. Sobre espectros sintéticos con punto de doblado
        fraccionario, el sesgo de la anchura recuperada baja del $+9{,}5\%$ con
        interpolación lineal al $+2{,}8\%$ con la cúbica. NORMOS trunca a canal
        entero y suma sin interpolar.
  \item \textbf{Dos ciclos de búsqueda.} El centro se localiza, se refina en una
        segunda pasada y solo entonces se dobla.
  \item \textbf{Los canales de borde se conservan.} Las versiones anteriores
        recortaban un número fijo de canales en cada extremo. Ya no se descartan
        por sistema: solo se detectan y recortan los canales extremos
        \emph{anómalos} --- los muertos o saturados que aparecen de vez en cuando
        en el primer o el último canal del multicanal.
\end{itemize}

\begin{truco}
Los canales de borde muertos no son una rareza. En un banco de trabajos reales
de NORMOS, 42 ajustes mejoraron hasta $\times 65$ al detectarlos, con el
$\chi^2$ reducido del peor caso cayendo de 102 a 1,45. Si un ajuste se niega a
converger a algo sensato, mire el primer y el último canal del espectro crudo
antes de culpar al modelo.
\end{truco}

\paragraph{Efecto geométrico.} La tasa de cuentas se modula con la
\emph{posición} del transductor, no solo con su velocidad. Fitbauer mide esa
modulación y la informa, igual que hace NORMOS bajo \emph{Geometry effect} en su
\ruta{.RES}, pero \textbf{no} la corrige: el efecto es antisimétrico respecto al
punto de doblado, así que el propio doblado lo cancela. Una amplitud relativa
por encima del $1\%$ delata un problema de geometría del espectrómetro y
conviene mirarlo antes de dar el ajuste por bueno.

\begin{aviso}
Un centro mal elegido produce líneas dobles o asimétricas. Si observa
desdoblamientos espurios o un espectro «fantasma» desplazado, revise
\textbf{primero} el centro de folding antes de tocar los parámetros del modelo.
\end{aviso}

\subsection{Forma de onda: triangular o senoidal}
\label{sec:drive-form}

Todo lo anterior asume un drive de \textbf{aceleración constante}, cuya onda de
velocidad es \textbf{triangular}: la velocidad sube y baja en rampas rectas, de
modo que avanza \emph{linealmente} con el canal y las dos mitades del barrido
recorren las mismas velocidades. Por eso tiene sentido doblar y usar un eje
lineal.

Algunos espectrómetros usan en cambio un drive \textbf{senoidal}. Entonces la
velocidad de cada canal no es lineal, sino
\[
  v_i \;=\; v_{\max}\,\sin\!\Big(\tfrac{2\pi\,(i-c_0)}{N}\Big),
\]
donde $N$ es el número de canales y $c_0$ la \textbf{fase} (el canal en que
$v=0$). Al ser el eje senoidal, doblar por simetría ya \emph{no} da un eje
lineal, y forzar uno lineal desplazaría las líneas.

\autofig[\textwidth]{fig_sine_drive}{Drive senoidal. Izquierda: la velocidad de
cada canal sigue una \emph{senoide}, no una rampa. Derecha: el espectro
resultante, con cada canal colocado en su velocidad real. No se pliega: se ajusta
tal cual.}

El control \menu{Forma de onda} del panel de calibración selecciona el modo:

\begin{itemize}[nosep]
  \item \textbf{Triangular} (por defecto): se dobla y el eje es lineal, como se ha
        descrito arriba.
  \item \textbf{Senoidal}: siguiendo el criterio de NORMOS (\emph{no doblar} y
        ajustar el espectro completo, con las dos ramas a la vez), Fitbauer asigna
        a cada canal su velocidad senoidal real y ajusta \textbf{sin plegar}.
\end{itemize}

Sobre la \textbf{fase} $c_0$: al igual que en triangular hay que localizar el
centro, en senoidal hay que localizar el punto de simetría del barrido. Fitbauer
lo autodetecta y lo expone en el control \param{Centro} (ajustable en vivo).
Como una fase mal puesta desplazaría todas las velocidades, conviene además
\textbf{refinarla en el ajuste}: al activar \emph{Ajustar centro}, $c_0$ entra
como \textbf{parámetro libre} y el optimizador la afina (igual que \emph{Ajustar
Vmax} afina la amplitud). Es lo que en NORMOS corresponde a \code{TRIANG=.FALSE.}
con \code{FOLD=.FALSE.} y \code{SIMULT=.TRUE.}

\begin{nota}
En modo senoidal el eje de velocidad no es monótono (la misma velocidad aparece
en varias fases del ciclo), por lo que los datos se dibujan como \textbf{puntos}
que recorren $-v_{\max}\!\dots\!+v_{\max}$ y vuelven. El modelo y el ajuste lo
tratan sin problema, porque evalúan el modelo \emph{punto a punto} (no necesitan
un eje ordenado). El componente de relajación no está soportado con drive
senoidal en esta versión.
\end{nota}

\section{Normalización y ruido}

Para que los ajustes sean comparables entre espectros, los datos se
\textbf{normalizan} dividiendo por un factor (el percentil 90 de las cuentas
dobladas), de modo que la línea base queda en $\approx 1$ y las absorciones se
leen como fracciones. Ese mismo factor se usa para estimar la incertidumbre de
cada punto a partir del ruido de Poisson (proporcional a $\sqrt{N}$), que después
pondera el ajuste (véase la \emph{verosimilitud}, §\ref{sec:opciones-ajuste}).

\begin{nota}
La \textbf{línea base} y la \textbf{pendiente} (control \param{Pendiente}) del
panel de calibración modelan una posible deriva del fondo. Se pueden fijar o
dejar libres en el ajuste como cualquier otro parámetro.
\end{nota}

\section{Visualización}

La zona de gráfico muestra:

\begin{itemize}[nosep]
  \item Los \textbf{datos} (puntos), la \textbf{línea base}, el \textbf{modelo}
        total y, cuando procede, los \textbf{subespectros} de cada componente
        (con su relleno de área opcional).
  \item Debajo, la \textbf{gráfica de diferencia} (residuo) datos$-$modelo.
  \item En modo distribución, además la \textbf{envolvente} de la distribución y
        cada componente nítido por separado, y un panel con $P(\BHF)$.
\end{itemize}

Desde el menú \menu{Vista} se controla la presentación:

\begin{itemize}[nosep]
  \item \menu{Vista > Mostrar diferencia} / \menu{Mostrar leyenda} /
        \menu{Mostrar área de componentes}: activan u ocultan cada elemento.
  \item \menu{Vista > Estilo de gráficos}: \emph{Clásico}, \emph{Moderno},
        \emph{Publicación} u \emph{Oscuro}.
  \item \menu{Vista > Tema visual} y \menu{Tema de color}: apariencia general de
        la interfaz.
  \item \menu{Vista > Configurar layout de paneles…}: reorganiza y guarda presets
        de disposición.
\end{itemize}

La \textbf{barra de herramientas} bajo el gráfico (Matplotlib) permite hacer
zoom, desplazarse, volver a la vista inicial y \textbf{guardar la figura} como
imagen (PNG, PDF, SVG).

\guicap{gui_plot_area}{Zona de gráfico con datos, modelo, subespectros y residuo,
con su barra de herramientas de navegación.}

\section{Comparar espectros}
\label{sec:comparar}

\menu{Archivo > Comparar espectro…} superpone uno o varios espectros de
referencia (con una paleta de colores propia) sobre el actual, útil para
identificar fases o comparar tratamientos térmicos. \menu{Archivo > Limpiar
comparación} retira todas las superposiciones. La comparación es solo visual: no
altera los datos ni el ajuste.

\section{Guardar, exportar y reutilizar}

\begin{center}
\begin{tabular}{@{}p{0.42\textwidth}p{0.5\textwidth}@{}}
\toprule
\textbf{Acción} & \textbf{Qué guarda} \\
\midrule
\menu{Archivo > Guardar sesión…} & Estado \textbf{completo} en \ruta{.json}
(datos, calibración fijada, modelo, opciones, distribución). Reabre con
\menu{Cargar sesión…}. \\[4pt]
\menu{Archivo > Guardar ajuste…} & El ajuste actual (parámetros, curva,
distribución) en texto. \\[4pt]
\menu{Archivo > Exportar informe Markdown/PDF…} & Informe legible con parámetros,
estadísticos y, si procede, la tabla de la distribución. \\[4pt]
\menu{Archivo > Exportar informe reducido…} & Versión abreviada del informe. \\
\bottomrule
\end{tabular}
\end{center}

\begin{truco}
La \textbf{sesión} es la vía recomendada para reproducibilidad: guárdela al
terminar un análisis y podrá retomarlo exactamente igual —incluida la
calibración fijada— en otro momento u ordenador. Si dispone de la API del
laboratorio, \menu{Archivo > Web} permite además descargar espectros/calibraciones
y \menu{Subir sesión JSON a web}.
\end{truco}

\section{Interoperabilidad con NORMOS: ficheros \texorpdfstring{\texttt{.JOB}}{.JOB}}
\label{sec:normos-job}

NORMOS es el programa con el que se ha analizado buena parte de la bibliografía
Mössbauer publicada, y muchos laboratorios guardan años de trabajo en ficheros
\ruta{.JOB}. Fitbauer lee y escribe ese formato. \textbf{No} ejecuta NORMOS ni
lo distribuye: solo habla su formato de texto, que no es propietario.

\subsection*{Importar}

\menu{Archivo > NORMOS (.JOB) > Importar trabajo de NORMOS…} lee un trabajo y
reconstruye el modelo en los paneles de Fitbauer. \textbf{El espectro se carga
también.} Un \ruta{.JOB} nombra sus ficheros en las cuatro primeras líneas, sin
ruta, porque NORMOS corría en DOS con todo en el mismo directorio; Fitbauer
busca ahí el espectro sin distinguir mayúsculas de minúsculas, y si el nombre
declarado no aparece prueba con el del propio trabajo y, en último caso, con el
único espectro de la carpeta. Deje todos los ficheros de un trabajo en la misma
carpeta y la importación funcionará de una vez.

Se reconocen solas dos familias de trabajo:

\begin{itemize}[nosep]
  \item \textbf{NORMOS-SITE} (sitios discretos). Cada subespectro se convierte en
        un componente --- singlete, doblete o sexteto --- con sus valores, sus
        casillas de \param{fijo} y sus ligaduras
        \code{NDEX}/\code{FACTOR}/\code{CONST}.
  \item \textbf{NORMOS-DIST} (distribuciones). Abren el panel
        $P(\BHF)$/$P(\dEQ)$ en vez del modo discreto. Sus \code{NSUB}
        \emph{no} son sitios: son los puntos de una \textbf{malla}. Fitbauer
        traduce la malla (origen y paso), la forma (histograma, gaussiana,
        binomial o fija), la correlación $\IS(x)$ y los anclajes de borde; los
        subespectros \emph{cristalinos} pasan a ser componentes nítidos.
\end{itemize}

\subsection*{Exportar}

\menu{Archivo > NORMOS (.JOB) > Exportar trabajo de NORMOS…} escribe el modelo
actual en formato NORMOS. Se ha comprobado que NORMOS acepta el fichero que
produce Fitbauer y reproduce la teoría original con diferencia exactamente cero.

\subsection*{Conversiones de convenio}

Es la parte delicada, y equivocarse aquí no da ningún error:

\begin{center}
\begin{tabular}{@{}p{0.26\textwidth}p{0.66\textwidth}@{}}
\toprule
\textbf{NORMOS} & \textbf{Significado en Fitbauer} \\
\midrule
\code{WID}, \code{W13}, \code{W23} & \code{WID} es la anchura de las líneas 3,4 y
\code{W13}/\code{W23} son relativas a ella; el $\Gamma_1$ de Fitbauer es la de
las líneas 1,6. La conversión es $\Gamma_1 = \code{WID}\cdot\code{W13}$. \\[4pt]
\code{D13}, \code{D23} & Razones de \textbf{área}. Los \param{int1}/\param{int2}
de Fitbauer son razones de \textbf{profundidad}. Coinciden solo si las anchuras
son iguales. \\[4pt]
\code{DEP} (o \code{ARE}) & El \textbf{área} del subespectro en mm/s, no una
profundidad. \\[4pt]
\code{NDEX}/\code{FACTOR}/\code{CONST} & Ligaduras en la numeración
\textbf{global} de NORMOS, $13 + 15\,(n-1)$. \\
\bottomrule
\end{tabular}
\end{center}

\paragraph{La escala de $\BHF$.} NORMOS deriva las posiciones de las líneas del
sexteto de los momentos nucleares; Fitbauer usa el patrón publicado de
$\alpha$-Fe. No difieren en un simple factor de escala. Para reproducir un
$\BHF$ de NORMOS exactamente, ajuste con el convenio de NORMOS activo
(\code{core.constants.sextet\_pattern("normos")}); la diferencia es de unos
$0{,}1$\,T.

\subsection*{El punto de doblado no se impone}

El \code{PFP} que trae un \ruta{.JOB} es la \textbf{semilla} de la búsqueda del
punto de doblado, no su resultado: NORMOS lo refina en dos ciclos, y en trabajos
reales acaba a más de un canal de lo que pedía el fichero. Fitbauer hace su
propia búsqueda --- que es el análogo correcto --- e informa del \code{PFP} solo
como dato. Escríbalo a mano en \param{Centro} si quiere forzarlo.

Hay una segunda sutileza. El punto refinado que NORMOS imprime en su \ruta{.RES}
tampoco es donde dobla: su rutina final trunca ese valor y suma canales enteros,
dejando el eje de simetría en $\lfloor \code{PFP} \rfloor + 0{,}5$. Tenerlo en
cuenta es lo que permite reproducir sus ajustes.

\subsection*{Lo que no se traslada}

El importador avisa de cada uno de estos puntos, porque lo que \emph{no} se ha
traducido importa tanto como lo que sí:

\begin{itemize}[nosep]
  \item Distribuciones Czjzek / Le Caër (\code{DISTRI=4}) y el modelo de vecinos
        de Billard--Chamberod (\code{METHOD=3}).
  \item Varios bloques de distribución solapados: Fitbauer maneja uno.
  \item El parámetro de suavizado \code{LAMDA}. El de NORMOS es absoluto y el
        $\alpha$ de Fitbauer adimensional, así que no hay conversión uno a uno:
        hay que fijarlo con la L-curve (sección~\ref{sec:lcurve}). La
        \emph{razón} \code{BETA}/\code{LAMDA} sí se conserva y es el anclaje de
        bordes.
  \item \code{DTQ} en las distribuciones de campo. NORMOS lo ignora ahí, así que
        trasladarlo metería una correlación que él nunca aplicó.
\end{itemize}

\begin{aviso}
Cuidado con los \ruta{.JOB} heredados. El formato de DIST no admite claves de
SITE como \code{NLINE}, \code{DEP}, \code{W13} o \code{W23}. Si vienen copiadas
de otro trabajo, NORMOS las lee y las descarta \textbf{sin decir nada}, de modo
que ese subespectro nunca entró en su ajuste. Fitbauer sí lo avisa.
\end{aviso}
